% ==============================================================================
% ИНСТРУКЦИИ ДЛЯ ИИ-ВЕРСТАЛЬЩИКА (AI TYPESETTING STRICT DIRECTIVES)
% ВНИМАНИЕ: СТРОГО ОБЯЗАТЕЛЬНО К ИСПОЛНЕНИЮ ПРИ ГЕНЕРАЦИИ ЛЮБОГО КОНТЕНТА!
% ==============================================================================
% 1. ПОЛНЫЙ ЗАПРЕТ MARKDOWN: Категорически запрещено использовать любые 
%    символы Markdown для форматирования. НИКАКИХ **, *, #, _, ` или >.
%    - Жирный текст: только \textbf{...}
%    - Курсив: только \textit{...}
%    - Заголовки: только \section{...}, \subsection{...}
%
% 2. СПИСКИ: Любые перечисления должны быть оформлены СТРОГО через окружения
%    \begin{itemize} ... \end{itemize} или \begin{enumerate} ... \end{enumerate}.
%
% 3. МАТЕМАТИЧЕСКАЯ ВЕРСТКА: Запрещено использовать обычный текст для формул. 
%    - Внутристрочные формулы: строго $...$
%    - Выключные формулы: строго \[ ... \] или \begin{equation}, \begin{align}.
%
% 4. СЕМАНТИЧЕСКАЯ РАЗМЕТКА: Текст лекций и семинаров должен быть разбит 
%    на логические блоки через теоремные окружения (уже стилизованные tcolorbox):
%    \begin{definition}, \begin{theorem}, \begin{lemma}, \begin{proof}, \begin{example}.
%
% 5. ТАБЛИЦЫ И МАТРИЦЫ: Игры в нормальной форме верстать через {NiceTabular} 
%    из пакета nicematrix с использованием \toprule, \midrule, \bottomrule.
% ==============================================================================

\documentclass[12pt, a4paper,openany, twoside]{book}

% --- Кодировки и локализация ---
\usepackage[T2A]{fontenc}
\usepackage[utf8]{inputenc}
\usepackage[english, main=russian]{babel} % ИСПРАВЛЕНО: Русский язык сделан основным

% --- Геометрия страницы и интерлиньяж ---
\usepackage[top = 2cm, bottom = 2cm, right =1cm, left =1cm, headheight=15pt]{geometry}
\usepackage{microtype}
\usepackage{setspace}
\onehalfspacing

% --- Математика ---
\usepackage{amsmath, amssymb, amsthm, mathtools, amsfonts, bm,mathrsfs}
\usepackage{mleftright} \mleftright
\usepackage{indentfirst} % ДОБАВЛЕНО: Красная строка для первого абзаца

% --- Графика и Таблицы ---
\usepackage{tikz}
\usepackage{istgame}
\usepackage{booktabs, multirow, array, tabularx, makecell}
\usepackage{nicematrix}
\usepackage{graphicx} % Для вставки фото

% --- Цвета и гиперссылки ---
\usepackage{xcolor}
\definecolor{darkblue}{RGB}{0, 51, 153}
\definecolor{darkred}{RGB}{153, 0, 0}
\definecolor{boxblue}{RGB}{240, 245, 255}
\definecolor{boxred}{RGB}{255, 240, 240}
\definecolor{boxgreen}{RGB}{240, 255, 240}
\definecolor{darkgray}{RGB}{105, 105, 105}
\definecolor{boxgray}{RGB}{248, 248, 248}
\definecolor{thmblue}{HTML}{2980B9}
\definecolor{ticketdark}{HTML}{2C3E50}

\usepackage{hyperref}
\hypersetup{
    unicode=true, % ДОБАВЛЕНО: Поддержка кириллицы в метаданных и закладках
    colorlinks=true,
    linkcolor=darkblue,
    citecolor=darkred,
    urlcolor=darkblue,
    filecolor=darkblue,
    pdftitle={Теория игр: Лекции и семинары},
    pdfpagemode=UseNone
}

% --- TCOLORBOX (Профессиональные рамки) ---
\usepackage[many]{tcolorbox}

% Настройка стандартных окружений amsthm
\theoremstyle{plain}
\newtheorem{theorem}{Теорема}[section]
\newtheorem{lemma}{Лемма}[section]

\theoremstyle{definition}
\newtheorem{definition}{Определение}[section]
\newtheorem{example}{Пример}[section]

\theoremstyle{remark}
\newtheorem{remark}{Замечание}[section]

% Привязка tcolorbox к окружениям amsthm (сохраняет стандартный синтаксис LaTeX)
\tcolorboxenvironment{definition}{
  enhanced, breakable,
  colback=boxblue, boxrule=0pt, frame hidden,
  borderline west={3pt}{0pt}{darkblue},
  sharp corners, before skip=10pt, after skip=10pt
}

\tcolorboxenvironment{theorem}{
  enhanced, breakable,
  colback=boxred, boxrule=0pt, frame hidden,
  borderline west={3pt}{0pt}{darkred},
  sharp corners, before skip=10pt, after skip=10pt
}

\tcolorboxenvironment{example}{
  enhanced, breakable,
  colback=boxgreen, boxrule=0pt, frame hidden,
  borderline west={3pt}{0pt}{green!50!black},
  sharp corners, before skip=10pt, after skip=10pt,
  % МАГИЯ ЗДЕСЬ: подменяем окружение table внутри примера
  before upper={\let\table\fixedtable\let\endtable\endfixedtable} 
}

\tcolorboxenvironment{proof}{
  enhanced, breakable,
  colback=boxgray, boxrule=0pt, frame hidden,
  borderline west={3pt}{0pt}{darkgray},
  sharp corners, before skip=10pt, after skip=10pt,
  % На всякий случай оставляем хак для таблиц, вдруг в доказательстве будет таблица
  before upper={\let\table\fixedtable\let\endtable\endfixedtable} 
}

% --- ХАКИ ДЛЯ СОВМЕСТИМОСТИ СО СТАРЫМ СИНТАКСИСОМ ---
\usepackage{caption} % Требуется для правильной работы подписей в неплавающих блоках

\makeatletter
% 1. Создаем статичную замену для плавающей таблицы
\newenvironment{fixedtable}[1][]{% игнорируем аргумент [h]
  \par\medskip
  \begin{center}%
  \def\@captype{table}% Подменяем тип для команды \caption
}{%
  \end{center}%
  \par\medskip
}

% 2. Учим nicematrix игнорировать ключ wd с помощью встроенных макросов LaTeX3
\ExplSyntaxOn
\keys_define:nn { NiceMatrix } { wd .code:n = {} }
\ExplSyntaxOff
\makeatother

% --- Пользовательские макросы ---
\DeclareMathOperator{\BR}{BR}
\newcommand{\strat}[1]{\mathcal{#1}}
\newcommand{\prof}[1]{\mathbf{#1}}

\setlength{\parindent}{1.25cm}
\setlength{\parskip}{0.5em}

\usepackage{import}

%Пример теха
%\begin{document}
%
%\section{Статические игры с полной информацией}
%
%В данном разделе мы рассмотрим базовые концепции теории игр на основе классических примеров. Важнейшей концепцией решения для игр в нормальной форме является равновесие, предложенное Джоном Нэшем.
%
%\begin{definition}[Равновесие Нэша]\label{def:nash}
%Пусть $G = \langle I, S, u \rangle$ — игра в нормальной форме. Тогда профиль стратегий $s^* \in S$ называется равновесием Нэша, если для всех игроков $i \in I$ и для всех альтернативных стратегий $s_i' \in S_i$ выполняется условие:
%\begin{equation}
%    u_i(s_i^*, s_{-i}^*) \ge u_i(s_i', s_{-i}^*)
%\end{equation}
%Иными словами, ни один игрок не может увеличить свой выигрыш, в одностороннем порядке отклонившись от равновесной стратегии.
%\end{definition}
%
%Как следует из Определения \ref{def:nash}, равновесие подразумевает отсутствие стимулов к отклонению. Рассмотрим классическую иллюстрацию этой концепции.
%
%\begin{example}[Дилемма заключенного]
%Два бандита, Петя и Вася, попались полиции. Они сидят в разных камерах и лишены возможности общаться друг с другом. У каждого есть выбор: «сознаться» или «молчать». Матрица выигрышей (где числа отражают полезность, отрицательные значения — сроки заключения) представлена в Таблице \ref{tab:prisoner}.
%
%\begin{table}[h]
%\centering
%\renewcommand{\arraystretch}{1.3}
%\begin{NiceTabular}{cc|cc}
%\toprule
%& & \multicolumn{2}{c}{\textbf{Вася}} \\
%& & \textit{«молчать»} & \textit{«сознаться»} \\
%\midrule
%\multirow{2}{*}{\textbf{Петя}} 
%& \textit{«молчать»} & $-1, -1$ & $-10, 0$ \\
%& \textit{«сознаться»} & $0, -10$ & $-8, -8$ \\
%\bottomrule
%\end{NiceTabular}
%\caption{Матрица выигрышей в игре «Дилемма заключенного»}
%\label{tab:prisoner}
%\end{table}
%
%Анализ матрицы показывает, что стратегия \textit{«сознаться»} строго доминирует стратегию \textit{«молчать»} для обоих игроков. Следовательно, единственным равновесием Нэша в данной игре является профиль (\textit{«сознаться»}, \textit{«сознаться»}) с выигрышами $(-8, -8)$, несмотря на то, что исход $(-1, -1)$ Парето-предпочтительнее.
%\end{example}
%
%\begin{theorem}[Связь доминирования и равновесия]
%Пусть в игре $G = \langle I, S, u \rangle$ профиль стратегий $s^* \in S$ является равновесием Нэша. Тогда $s^*$ выживает в процессе итеративного удаления строго доминируемых стратегий.
%\end{theorem}
%
%\begin{proof}
%Предположим противное: пусть $s^*$ — равновесие Нэша, но оно удаляется на шаге $n$. Тогда существует игрок $i$ и стратегия $s_i' \in S_i^n$ такая, что для всех $s_{-i} \in S_{-i}^n$ выполняется $u_i(s_i', s_{-i}) > u_i(s_i^*, s_{-i})$. 
%Однако это противоречит определению равновесия Нэша, согласно которому $u_i(s_i^*, s_{-i}^*) \ge u_i(s_i', s_{-i}^*)$. Следовательно, равновесный профиль не может быть удален.
%\end{proof}
%
%\end{document}

\begin{document}

\begin{titlepage}
    % Временно убираем поля для обложки, чтобы залить ее цветом на 100%
    \newgeometry{left=0cm, right=0cm, top=0cm, bottom=0cm}
    
    \begin{tikzpicture}[remember picture, overlay]
        % 1. Вставляем фотографию на верхние 60% страницы
        % Scope используется для того, чтобы аккуратно обрезать фото, если оно не влезает, без искажения пропорций
        \begin{scope}
            \clip (current page.north west) rectangle ([yshift=-0.6\paperheight]current page.north east);
            \node[anchor=north, inner sep=0pt] at (current page.north) {
                \includegraphics[width=\paperwidth]{photo.png} % <-- ВСТАВЬ ИМЯ СВОЕГО ФАЙЛА ЗДЕСЬ
            };
        \end{scope}
        
        % 2. Темный блок на нижние 40% страницы (цвет ticketdark)
        \fill[ticketdark] (current page.south west) rectangle ([yshift=-0.6\paperheight]current page.north east);
        
        % 3. Тонкая акцентная полоса (thmblue) между фото и текстом для эстетики
        \fill[thmblue] ([yshift=-0.6\paperheight]current page.north west) rectangle ([yshift=-0.6\paperheight-3pt]current page.north east);

        % 4. Типографика обложки (выровнена по левому краю)
        \node[anchor=north west, text=white, align=left] at ([xshift=2.5cm, yshift=-0.68\paperheight]current page.north west) {
            {\sffamily\Huge\bfseries ТЕОРИЯ ИГР}\\[0.8em]
            {\sffamily\Large Конспект лекций теории игр}\\[2.5em]
            {\sffamily\large\textcolor{thmblue!50}{Автор:}} {\sffamily\large Dr\_papayka}\\[0.3em]
        };
    \end{tikzpicture}
    
    % Возвращаем стандартные поля для остального документа
    \restoregeometry
\end{titlepage}


\tableofcontents
\newpage

\chapter{1 лекция}
\import{conspect}{1.tex}
\chapter{2 лекция}
\import{conspect}{2.tex}
\chapter{3 лекция}
\import{conspect}{3.tex}
\chapter{4 лекция}
\import{conspect}{4.tex}
\chapter{5 лекция}
\import{conspect}{5.tex}
\chapter{6 лекция}
\import{conspect}{6.tex}

\chapter{7 лекция}
\import{conspect}{7.tex}

\chapter{8 лекция}
\import{conspect}{8.tex}

\chapter{9 лекция}
\import{conspect}{9.tex}



\end{document}
